% Procedencia: integral 2249, cap. 25, pp. 619--640, especialmente §25.13.
% Requiere la incorporación material de exc:k-direccion, exc:leech y el núcleo común.
\section{Écrans dimensionnels et réduction coordonnée de l’état}
\label{iv:pantallas}

Les rangs qui apparaissent dans le prolongement exceptionnel et dans la
compactification appartiennent à des objets différents. Pour les relier,
on conserve leur présentation : générateurs, relations, produit
scalaire et direction sélectionnée. La dimension est alors une
conséquence du module des canaux et de ses relations. Son
transport est formulé au moyen d’applications qui agissent sur ce module,
en plus d’agir sur l’état qui le détermine.

\subsection{Canaux, relations et dimension effective}

Soit \(\mathcal S_n\) l’ensemble des états survivants de
profondeur \(n\), avec les restrictions compatibles
\(r_{m,n}:\mathcal S_m\to\mathcal S_n\), \(m\geq n\).
La donnée de frontière \(\beta_n(x)\) conserve l’information
nécessaire pour décider de la compatibilité d’un prolongement :
feuillet, orientation, quotient, mémoire et classe d’incidence selon
le lecteur considéré. Les restrictions vérifient
\[
 r_{\ell,n}=r_{m,n}r_{\ell,m},\qquad
 \beta_n r_{m,n}=b_{m,n}\beta_m.
\]
Ces égalités sont les conditions de naturalité de la
présentation et maintiennent les mêmes données lorsqu’on restreint par
étapes ou directement.


\begin{definition}[Écran de canaux]
\label{iv:def:pantalla}
Un corps \(\mathbb K_n\) étant fixé, soit \(M_n(x)\) l’espace
engendré par les canaux de prolongement et \(\mathcal R_n(x)\)
le sous-espace des relations. L’écran est
\[
 \Sigma_n(x)=\bigl(\mathbb K_n,M_n(x),\beta_n(x),
                    \mathcal R_n(x)\bigr),
\]
et sa dimension effective est
\begin{equation}
 d_{\Sigma_n}(x)=\dim_{\mathbb K_n}
                    \bigl(M_n(x)/\mathcal R_n(x)\bigr).
 \label{iv:eq:dimension-pantalla}
\end{equation}
Si \(M_n(x)\cong\mathbb K_n^{g_n}\) et que les colonnes de
\(R_n(x)\) engendrent les relations, alors
\(d_{\Sigma_n}(x)=g_n-\operatorname{rank}R_n(x)\).
\end{definition}

\begin{proposition}[Invariance de présentation et transport
]
\label{iv:prop:pantalla-transporte}
La dimension \eqref{iv:eq:dimension-pantalla} est conservée
sous des changements de bases inversibles. Pour deux écrans sur un même corps
\(\mathbb K_n=\mathbb K_m=\mathbb K\), si l’application linéaire
\(F_\gamma:M_n(x)\to M_m(T_\gamma x)\) envoie
\(\mathcal R_n(x)\) dans \(\mathcal R_m(T_\gamma x)\), elle induit
de manière unique une application linéaire entre les quotients d’écran.
\end{proposition}
\begin{proof}
Un changement de générateurs et de relations transforme \(R\) en
\(URV\), avec \(U,V\) inversibles, et conserve son rang. Pour
la descente, nous définissons
\(\overline F_\gamma([v])=[F_\gamma(v)]\). Si deux
représentants diffèrent d’une relation, la linéarité implique que leurs images diffèrent
d’une relation de l’écran terminal ; l’application est
bien définie. La projection sur le quotient est surjective et
garantit l’unicité.

\end{proof}

Le changement dimensionnel pendant le transport est déterminé
par
\[
 \Delta d(\gamma;x)=d_{\Sigma_m}(T_\gamma x)
                         -d_{\Sigma_n}(x).
\]
Cette expression conserve l’état et sa présentation de canaux.
Un rang isolé ne permet de reconstruire ni les relations ni
l’application qui l’a produit.

\subsection{Poinçonnage ternaire et écran de syndromes}

Le code \(C_W=[12,6,6]_3\), construit à partir de
l’incidence exceptionnelle, produit un autre écran par
poinçonnage d’une coordonnée. Notons
\(p:C_W\to\FF_3^{11}\) cette application et
\(G_{11}=p(C_W)\) son image.

\begin{proposition}[Code poinçonné et partition parfaite]
\label{iv:prop:codigo-once}
L’application \(p\) est injective. Le code \(G_{11}\) a pour
paramètres \([11,6,5]_3\), et ses boules de Hamming de rayon
\(2\) forment une partition de \(\FF_3^{11}\).
\end{proposition}
\begin{proof}
Un vecteur du noyau de \(p\) aurait son support contenu dans
l’unique coordonnée retirée. La distance minimale \(6\)
de \(C_W\) exclut tout vecteur non nul de cette forme.
Ainsi, la dimension reste \(6\). Le poinçonnage diminue
le poids d’au plus une unité, donc la distance minimale
de l’image vaut au moins \(5\). Les boules de rayon \(2\)
sont donc disjointes. Chacune contient

\[
 1+2\binom{11}{1}+2^2\binom{11}{2}
 =1+22+220=243
\]
vecteurs. Comme \(|G_{11}|=3^6\) et
\(3^6\cdot243=3^{11}\), les boules recouvrent l’espace ambiant.
Pour vérifier que la distance est exactement \(5\),
prenons un vecteur de poids \(3\). Sa boule de décodage
ne peut avoir pour centre zéro et possède un centre de poids au
plus \(3+2=5\) ; ce centre doit avoir un poids au moins égal à
\(5\). On obtient ainsi un mot de poids \(5\).
\end{proof}

L’entier \(243=3^{11-6}\) est ici le nombre de classes
du quotient des syndromes \(\FF_3^{11}/G_{11}\). Sa
fonction est fixée par le poinçonnage et la distance du
code ; l’égalité de cardinal avec une autre fibre ne remplace pas
l’application entre les deux. Cet écran combinatoire et
l’écran réel suivant conservent leurs opérations propres.

\subsection{La direction dodécaphasique et la réduction du rang}

La construction de \ref{exc:k-direccion} part du registre
\(K\) déjà obtenu par APP--TRIT--TPK et du système de coordonnées de
représentation déclaré sur \(\RR^{12}\). Dans ce système de coordonnées,
\begin{equation}
 V_{11}=\mathbf1^\perp,\qquad
 P_{11}=I_{12}-\frac1{12}J_{12},\qquad
 u_K=P_3P_{11}K.
 \label{iv:eq:p11-uk}
\end{equation}
Ici, \(J_{12}\) est la matrice de uns et \(P_3\) est le
projecteur isotypique spécifié par l’action de \(A_5\).
La somme finie qui définit \(P_3\), ses représentants et la
génération de \(K\) sont conservés dans la section citée de
ce même article. En particulier,

\[
 \|u_K\|^2=\frac{6638585+2275584\sqrt5}{20}>0.
\]
La norme positive définit \(\ell_K=\RR u_K\), et le
projecteur
\begin{equation}
 P_{10}=P_{11}-\frac{u_Ku_K^{\mathsf T}}{\|u_K\|^2}
 \label{iv:eq:p10}
\end{equation}
a pour image \(V_{10}=V_{11}\cap u_K^\perp\).

\begin{proposition}[Deux réductions orthogonales]
\label{iv:prop:12-11-10}
On a
\[
 \RR^{12}=\RR\mathbf1\mathbin{\oplus^\perp}V_{11},\qquad
 V_{11}=\ell_K\mathbin{\oplus^\perp}V_{10},\qquad
 \dim V_{10}=10.
\]
De plus, \(P_{10}P_{11}=P_{10}\), \(P_{10}u_K=0\), et
\(P_{10}\) est un projecteur orthogonal.
\end{proposition}
\begin{proof}
La proposition \ref{prop:k-reduccion-diez}, démontrée avant de
construire les écrans, établit l’idempotence, la symétrie, le rang et
les deux identités opératorielles pour ce même \(u_K\) et ce même
\(P_{10}\). La première décomposition s’obtient à partir de
\(P_{11}=I-J_{12}/12\) ; la seconde est la décomposition orthogonale
de cette proposition. On conserve ainsi la démonstration unique de la
réduction et on explicite ici ses deux étapes.
\end{proof}

La première réduction élimine exclusivement le mode uniforme.
La seconde élimine la direction déterminée par \(K\) dans
la représentation fixée. Si l’on change simultanément le système de coordonnées
par une matrice orthogonale \(S\), le vecteur se transforme par
\(u_K\mapsto Su_K\) et le projecteur par
\(P_{10}\mapsto SP_{10}S^{-1}\). Cette covariance conserve
l’opération géométrique, même si ses coordonnées changent.

\subsection{Écran réticulaire lorentzien
}

Soit \(\Lambda=N_\alpha\) le réseau de Leech construit
dans \ref{exc:leech}. On considère le plan hyperbolique entier
\(U=\ZZ e_+\oplus\ZZ e_-\), avec

\[
 \langle e_+,e_+\rangle=\langle e_-,e_-\rangle=0,
 \qquad\langle e_+,e_-\rangle=1.
\]
Le réseau \(L_{26}=\Lambda\oplus U\) est pair et
unimodulaire, de signature \((25,1)\). Chaque vecteur s’écrit
\(y=\lambda+ae_++be_-\). Comme
\(\langle y,e_+\rangle=b\),

\begin{equation}
 e_+^\perp=\Lambda\oplus\ZZ e_+,
 \qquad e_+^\perp/\ZZ e_+\cong\Lambda.
 \label{iv:eq:26-24}
\end{equation}
L’application du quotient est
\(q_{24}(\lambda+ae_+)=\lambda\) ; elle est surjective et a pour
noyau \(\ZZ e_+\). La diminution du rang de deux provient
d’une restriction d’orthogonalité suivie d’un quotient
isotrope. C’est donc une opération de type différent de
\eqref{iv:eq:p10}.

\begingroup\fontsize{10}{11.7}\selectfont
\setlength{\parskip}{2pt}\setlength{\abovedisplayskip}{5pt}\setlength{\belowdisplayskip}{5pt}
\subsection{Morphisme conjoint d’écrans
}

Avec le domaine \(\mathscr D_0\) de \ref{iv:def:dualidad},
soient
\[
 \mathscr G_K=\{(v_{11},P_{10}v_{11}):v_{11}\in V_{11}\},
\]
\[
 \mathscr S_{\mathrm{HMT}}=\RR^{12}\times\mathscr D_0\times e_+^\perp,
 \qquad
 \mathscr S_M^{\mathrm{alg}}=\mathscr G_K\times\mathscr D_0\times\Lambda.
\]
\begin{definition}[Application algébrique d’écrans coordonnés
]
\label{iv:def:rm-alg}
On définit
\begin{equation}
 \mathcal R_M^{\mathrm{alg}}(v,d,y)
 =\bigl((P_{11}v,P_{10}v),d,q_{24}(y)\bigr).
 \label{iv:eq:rm-alg}
\end{equation}
\end{definition}

\begin{theorem}[Isomorphisme quotienté et entrelacement de dualité
]
\label{iv:thm:pantallas-isomorfas}
\label{thm:k-pantallas-coordinadas}
L’application \eqref{iv:eq:rm-alg} induit un isomorphisme
\[
 \overline{\mathcal R}_M^{\mathrm{alg}}:
 (\RR^{12}/\RR\mathbf1)\times\mathscr D_0
 \times(e_+^\perp/\ZZ e_+)
 \longrightarrow\mathscr S_M^{\mathrm{alg}}.
\]
Il entrelace les transformations qui agissent par \(\mathsf T_0\)
sur le facteur \(\mathscr D_0\) et par l’identité sur les
deux autres facteurs. La forme \(E_{R_0}(m,w)\) est conservée
sous cette dualité.

\end{theorem}
\begin{proof}
L’application \([v]\mapsto P_{11}v\) est un isomorphisme
de \(\RR^{12}/\RR\mathbf1\) sur \(V_{11}\).
Comme \(P_{10}P_{11}=P_{10}\), l’application
\([v]\mapsto(P_{11}v,P_{10}v)\) identifie ce quotient
à \(\mathscr G_K\) ; son inverse prend la première composante
et sa classe. Par \eqref{iv:eq:26-24}, \(q_{24}\) induit
l’isomorphisme réticulaire du troisième facteur. Le deuxième est
transporté au moyen de l’identité. Le produit des trois
applications prouve l’isomorphisme.

Les projecteurs et \(q_{24}\) ne modifient pas \(d\), tandis que
la dualité agit uniquement sur \(d\). Les deux
compositions sont donc identiques. La conservation de la forme
est le théorème \ref{iv:thm:dualidad-escalar}.
\end{proof}

Le graphe \(\mathscr G_K\) conserve la coordonnée précédente
\(v_{11}\) avec sa projection. Le théorème démontre
l’équivalence des présentations quotientées ; la réduction
isolée \(V_{11}\to V_{10}\) conserve son noyau
\(\ell_K\). Cette précision établit quelles informations
retient l’isomorphisme et quelles informations élimine chaque
projection qui le compose.


\begin{corollary}[Écrans sur la collection elliptique]
Le même morphisme se restreint en un isomorphisme
\[
 (\RR^{12}/\RR\mathbf1)\times\mathscr D_{\mathrm{el}}
 \times(e_+^\perp/\ZZ e_+)
 \ \cong\ \mathscr G_K\times\mathscr D_{\mathrm{el}}\times\Lambda,
\]
et conserve l’entrelacement de dualité.
\end{corollary}
\begin{proof}
Le morphisme et son inverse conservent exactement le deuxième facteur.
Le corollaire \ref{iv:cor:radio-especializado} prouve que la dualité
préserve \(\mathscr D_{\mathrm{el}}\). On restreint donc
aussi bien la bijection que le carré d’entrelacement déjà démontré.
\end{proof}

\subsection{Coordination avec la réalisation exceptionnelle}

Sur le domaine \(\mathcal X_{\mathrm{HMT}}\), avec phase, orientation,
quotient, bord, incidence, graduation et mémoire, les lectures
\(\varepsilon_{\mathrm{exc}}\) et \(\varepsilon_{\mathrm{scr}}\) extraient
les données exceptionnelles et d’écran. Le corridor construit l’algèbre
graduée \(\mathfrak V(x)\cong V^\natural\), avec son produit de vertex
et ses automorphismes ; \(\mathfrak S_M(x)\) réunit les écrans, les quotients
et les opérateurs précédents. L’attribution conjointe
\(x\mapsto(\mathfrak V(x),\mathfrak S_M(x))\) a donc une
algèbre pour première composante et une collection d’espaces et d’applications pour seconde.
L’action du Monstre appartient à la première ; la réduction dirigée par \(K\)
et la dualité compacte, à la seconde. Les deux lectures fixent la provenance
commune et chaque construction conserve son type.
Le théorème~\ref{iv:thm:pantallas-isomorfas} donne l’isomorphisme
facteur par facteur ; sa preuve identifie les inverses, les noyaux
et les opérations conservées par chaque composante.

\par\endgroup
